\scanpage{23}
\mspage{\unsure{4|9}}
% archive foliation: 22

\uline{Mais attention} Il n'est pas vrai \uline{dans $\mathrm{Cat}$} que si on a un diagramme
cocartésien
\diagram{\begin{tikzcd}
Y\arrow[r,hook,dash,"i"]\arrow[d,"g"']&X\arrow[d,dash,"f"]\\
Y'\arrow[r,hook,dash,"i'"']&X'
\end{tikzcd}}
$i$ imm. fermée (ou ouverte), alors
\[
[\ i\in\underline W\ \Longrightarrow\ i'\in\underline W,\quad\text{ni que}\ ]\quad
 g\in\underline W\ \Longrightarrow\ f\in\underline W,
\]
$=$ pour l'\unsure{équiv.} faible habituelle, et $=$
si $g$ est \deleted{\unsure{fibré}}\inserted{cofibré}, \deleted{\unsure{fibré}}
et $i$ \unsure{une triv.} imm. fermée, $Y'\simeq\Delta^0$.

\margin{Si $i\in\underline W$, a-t-on pourtant $i'\in\underline W$ ?
(\uline{Ex} $Y,X$ asphériques $Y'=\Delta_0$) \uuline{Non}}
\margin{Il semble vrai pourtant que si $g$ induit une \deleted{iso} $1$-équiv.,
i.e. équiv. sur les groupoïdes fondamentaux, alors $f$ aussi. $[\ ]$ ?}

\uline{Ex} Considérons une $Z$ dans $(\mathrm{Cat})$, et
\diagram{\begin{tikzcd}[column sep=large,row sep=large,execute at end picture={\draw ([xshift=-4mm]\tikzcdmatrixname-2-1.north west) .. controls ([xshift=-8mm]\tikzcdmatrixname-2-1.north west) and ([xshift=-4mm]\tikzcdmatrixname-2-1.south west) .. ([xshift=-8mm]\tikzcdmatrixname-2-1.south west) .. controls ([xshift=-4mm]\tikzcdmatrixname-3-1.north west) and ([xshift=-8mm]\tikzcdmatrixname-3-1.south west) .. ([xshift=-4mm]\tikzcdmatrixname-3-1.south west);}]
Z\arrow[r,hook,dash,"\alpha", "{\text{imm. fermée}}"']\arrow[d,"\alpha"']\arrow[dr,phantom,"{\text{cocart}}" description]
&\check C(Z)\arrow[d]\arrow[r]&[\deleted{\unsure{e}}\ e_{u'_0}]\\
Y=\check C(Z)\arrow[r,hook,"{\text{imm. f. (\unsure{surj.})}}"']\arrow[d,"{g\in W}"']\arrow[dr,phantom,"{\text{cocart}}" description]
&\unsure{\overline\Sigma'}(Z)=X\arrow[d,"f"]&[e_{w_0}]\arrow[l,hook]\\
Y'=e\arrow[r,hook,dash,"{\text{imm. f.}}"']&\unsure{\Theta(Y')}=X'&[e_{u'_0}]\arrow[l,hook]

\end{tikzcd}}


Ici $g\in W$, et $X'\simeq\Delta^1$ ssi \deleted{\unsure{\pi_0}}
$Z$ \deleted{\unsure{n'}}est $0$-connexe.
(En général, $\mathrm{Hom}(e_{u'_0},e_{y'_0})\simeq\Pi_0(Z)$).
Mais \unsure{quand il en est ainsi}, $f\in\underline W$ ssi
$\unsure{\overline\Sigma}(Z)$ est $\underline W$-asph., ce qui n'est pas tjrs le cas.
(prendre p.ex. \deleted{\unsure{et}} $Z$ un groupoïde réduit à un \unsure{sglt.|sgl.}, avec $\Pi_1\ne1$)
